\documentclass[11pt]{article}
\usepackage[letterpaper,margin=1in]{geometry}
\usepackage[T1]{fontenc}
\usepackage{lmodern}
\usepackage{microtype}

\setlength{\parindent}{0pt}
\setlength{\parskip}{6pt}

\begin{document}

\textbf{Title:} ATHENA-rods: A Low-Cost Cyber-Physical Platform for Nuclear Control Education and Secure Digital I\&C Prototyping

\textbf{Authors:} Ondrej Chvala

\textbf{Email:} ochvala@utexas.edu

\textbf{Key findings:} A home-fabricable, approximately \$390 platform built from two Raspberry Pi 5 nodes and 3D-printed hardware can integrate reactor kinetics, physical actuation, and X.509-based secure communications for interdisciplinary engineering training.

Digital instrumentation and control (I\&C) modernization in nuclear systems requires engineers who can work across reactor dynamics, hardware integration, and cybersecurity. At the same time, there is a persistent gap between modern network-security capabilities and commonly deployed SCADA/OT practice, and educational infrastructure that makes this gap tangible is limited. ATHENA-rods addresses this need with an open-source cyber-physical platform for interdisciplinary training. This work is part of the ongoing nuclear digital twins effort at The University of Texas at Austin.

The platform is organized as three cooperating nodes: an instrument node (\texttt{instbox}), a control node (\texttt{ctrlbox}), and a visualization node (\texttt{visbox}). The instrument node uses a DC motor, servo, ultrasonic ranging, DHT11 environmental sensing, a limit switch, and an LED matrix to emulate control-rod motion and reactor-state feedback via real-time point kinetics. The control node manages authorization and command routing; the visualization node provides live telemetry and supervisory controls. Security is integrated through X.509 certificate management and TLS/mTLS communication, with optional face-plus-RFID operator authorization. Mechanical components are distributed as FreeCAD and STL assets for consumer FDM printing.

Current validation includes reproducibility manifests, control-performance metrics, and security verification workflows. Representative step-response tests show stable positive and negative reactivity transients with overshoot of approximately 9.0\% and 4.0\%, respectively. Estimated build cost is approximately \$390. Ongoing work includes hardware-backed latency/throughput benchmarking, certificate lifecycle testing, and DNP3/TLS integration. Code, CAD, and test artifacts are available on GitHub (https://github.com/ondrejch/athena-rods).

\end{document}
